\section{Conclusion}

The locally connected waveform correction improves the blind-test coincidence timing resolution (CTR) for both front-end boards and for both waveform families. The largest gains are obtained from the energy channel. With the wide waveform window, the relative CTR improvement ranges from approximately \(23\%\) to \(34\%\) for the UC board and from \(18\%\) to \(21\%\) for the FBK board across the investigated bias voltages. For example, the UC-board CTR at \SI{48}{V} is reduced from \(166\pm2\) ps with leading-edge discrimination (LED) to \(128\pm2\) ps after correction, while the FBK-board CTR at the same bias voltage is reduced from \(144\pm2\) ps to \(114\pm2\) ps.

Most of the energy-channel correction is already obtained from the short Onishi window around the leading edge. Extending the input to the wide window produces a further reduction of only \(1\)--\(3\) ps for the UC datasets and \(3\)--\(7\) ps for the FBK datasets. This behavior is consistent with the XAI maps, in which the strongest energy-channel sensitivity is concentrated around the rising edge. The threshold-scan studies also show that the learned correction substantially reduces the dependence on the particular LED threshold. In the representative UC energy-channel scan, the uncorrected LED CTR spans \(166\)--\(264\) ps over the tested thresholds, whereas the corrected CTR remains within \(129\)--\(141\) ps. For the corresponding FBK scan, the ranges are \(143\)--\(211\) ps before correction and \(114\)--\(121\) ps after correction.

The timing-channel case is different. Because this channel is specifically designed to provide a fast and relatively uniform timing signal, the correction obtainable from the short waveform region is smaller: the Onishi-window improvement is approximately \(1.0\)--\(1.8\%\) for the UC board and \(1.5\)--\(4.0\%\) for the FBK board. However, extending the waveform window consistently increases the improvement to approximately \(3.0\)--\(5.3\%\) for UC and \(5.1\)--\(6.5\%\) for FBK. The wide-window correction reaches \(57.4\pm0.8\) ps for the UC board at \SI{48}{V} and \(58\pm1\) ps for the FBK board at \SI{49}{V}. The representative XAI maps show sensitivity not only around the leading edge but also in later portions of the timing waveform, extending to approximately \SI{20}{ns} after the aligned crossing. This suggests that waveform structures following the primary rising edge, including secondary pulse features, may contain information correlated with the event-by-event timing error.

Overall, the results support the use of a shared antisymmetric LC-MLP as a lightweight waveform-based correction model and show that the useful temporal information depends strongly on the waveform family. The energy channel is dominated by information close to the leading edge, whereas the timing channel appears to benefit from additional temporal context. The planned comparison with the Onishi CNN under identical input windows and evaluation conditions will help determine whether these observations are specific to the locally connected architecture or reflect more general properties of the detector waveforms.
